% Continuation: R4, R5, references, appendix (included by report.tex)

\section{R4: Evaluation --- Non-Collaborative Baseline}

\subsection{Evaluation Metrics}

The following metrics are used to evaluate agent performance:

\begin{table}[htbp]
  \centering
  \caption{Evaluation metrics}
  \begin{tabular}{lp{0.62\linewidth}}
    \toprule
    \sustechheaderrow
    \sustechtableheader{Metric} & \sustechtableheader{Description} \\
    \midrule
    \textbf{Completion Time} & Total ticks from simulation start until all holes are filled (or max ticks reached) \\
    \textbf{Total Score} & Sum of rewards from all filled holes \\
    \textbf{Reward Efficiency} & Total score divided by completion ticks (score per tick) \\
    \textbf{Safety Rate} & Percentage of agents that never run out of energy during the simulation \\
    \textbf{Path Efficiency} & Ratio of actual distance traveled to theoretical shortest path \\
    \bottomrule
  \end{tabular}
\end{table}

\subsection{Single-Expectimax Utility (SEU) Function}

Following the original Tileworld formulation (Pollack \& Ringuette, 1990), the reward-based task selection strategy uses a \textbf{Single-Expectimax Utility (SEU)} function to evaluate and rank potential tasks. The SEU function combines expected reward with expected cost:

\begin{verbatim}
SEU(task) = E[Reward] - ExpectedCost
\end{verbatim}

Where the utility calculation incorporates:

\begin{enumerate}
  \item \textbf{Reward Component}: The hole's score value (1--10 points)
  \item \textbf{Travel Cost}: Estimated distance to complete the task chain (acquire tile $\to$ fill hole)
  \item \textbf{Energy Risk}: Post-task energy margin after completing the task and reaching nearest station
  \item \textbf{Switching Penalty}: Penalty for abandoning current target to pursue a new one (prevents thrashing)
\end{enumerate}

The implementation in \texttt{ContractNetBoard} uses a weighted formula:

\begin{verbatim}
Utility = w_score x norm(score)
        + w_hold  x (already_holding_tile ? 1 : 0)
        - w_dist  x norm(distance)
        - w_risk  x risk(slack)
        - w_switch x switchPenalty
\end{verbatim}

This SEU-based selection is compared against a \textbf{First-Available} baseline that selects the first unclaimed hole encountered, regardless of distance, reward, or energy feasibility.

\subsection{Experimental Setup}

\subsubsection{Parameter Configurations}

The implementation provides four predefined environment profiles via \texttt{EnvironmentConfig.java}:

\begin{table}[htbp]
  \centering
  \footnotesize
  \caption{Predefined environment profiles}
  \begin{tabular}{lcccc}
    \toprule
    \sustechheaderrow
    \sustechtableheader{Parameter} & \sustechtableheader{BASELINE} & \sustechtableheader{FAST\_CHANGE} & \sustechtableheader{LONG\_RANGE} & \sustechtableheader{CUSTOM} \\
    \midrule
    Grid Size & $100{\times}100$ & $50{\times}50$ & $150{\times}40$ & $100{\times}100$ \\
    Wind Interval & 20 ticks & 10 ticks & 40 ticks & 20 ticks \\
    Max Ticks & 5000 & 5000 & 5000 & 5000 \\
    Agent Sensing Range & 15 & 10 & 18 & 5 \\
    Energy Stations & 8 & 4 & 9 & 1 \\
    Agent Max Energy & 260 & 170 & 320 & 100 \\
    Low Energy Threshold & 110 & 65 & 140 & 20 \\
    Random Seed & 999 & 42 & 123 & 42 \\
    \bottomrule
  \end{tabular}
\end{table}

For R4 evaluation, we use the \textbf{BASELINE} profile as the default configuration and vary specific parameters to test sensitivity:

\begin{table}[htbp]
  \centering
  \caption{R4 sensitivity parameters}
  \begin{tabular}{lcc}
    \toprule
    \sustechheaderrow
    \sustechtableheader{Parameter} & \sustechtableheader{Baseline Value} & \sustechtableheader{Variations} \\
    \midrule
    Number of Agents & 5 & 3, 7, 10 \\
    Number of Holes & 10 & 5, 15, 20 \\
    Number of Obstacles & 15 & 5, 25, 40 \\
    Number of Tiles & Equal to holes & --- \\
    Wind Interval & 20 & 10 (fast), 50 (slow), $\infty$ (static) \\
    Max Ticks & 5000 & --- \\
    \bottomrule
  \end{tabular}
\end{table}

\subsubsection{Experimental Procedure}

\begin{itemize}
  \item 10 independent runs per configuration
  \item All results reported as average $\pm$ standard deviation
  \item Contract Net coordination disabled for baseline evaluation (R4)
  \item Random seed varied across runs for statistical robustness
\end{itemize}

\subsection{Results: Reward-Based vs First-Available Selection}

\tileworldfigure{Completion time comparison between First-Available and Reward-Based strategies across different configurations. Lower is better; * indicates the winner.}{assets/r4_completion_time.png}{fig:r4-time}

\begin{table}[htbp]
  \centering
  \caption{Average completion time (ticks)}
  \begin{tabular}{p{0.34\linewidth}ccc}
    \toprule
    \sustechheaderrow
    \sustechtableheader{Configuration} & \sustechtableheader{First-Avail.} & \sustechtableheader{Reward-Based} & \sustechtableheader{Winner} \\
    \midrule
    Baseline ($50{\times}50$, 5 agents, 15 obstacles) & 1309.1 & 1450.9 & First-Available \\
    Low obstacles (5) & 844.5 & 538.3 & \textbf{Reward-Based} \\
    High obstacles (40) & 974.1 & 376.0 & \textbf{Reward-Based} \\
    Fast wind (10 ticks) & 1059.2 & 1424.3 & First-Available \\
    Static environment & 5000.0 (4.0 holes left) & 5000.0 (2.4 holes left) & Reward-Based (more complete) \\
    3 agents & 964.9 & 1364.9 & First-Available \\
    10 agents & 550.7 & 347.2 & \textbf{Reward-Based} \\
    \bottomrule
  \end{tabular}
\end{table}

\textbf{Key finding}: The advantage of reward-based selection is \textbf{context-dependent}: high obstacle density and more agents favor Reward-Based; baseline and fewer agents favor First-Available.

\tileworldfigure{Reward efficiency comparison between First-Available and Reward-Based strategies. Higher is better; * indicates the winner.}{assets/r4_reward_efficiency.png}{fig:r4-eff}

\begin{table}[htbp]
  \centering
  \caption{Average reward efficiency (score/tick)}
  \begin{tabular}{p{0.28\linewidth}ccc}
    \toprule
    \sustechheaderrow
    \sustechtableheader{Configuration} & \sustechtableheader{First-Avail.} & \sustechtableheader{Reward-Based} & \sustechtableheader{Winner} \\
    \midrule
    Baseline & 39.58 & 24.97 & First-Available \\
    Low obstacles & 34.76 & 31.82 & First-Available \\
    High obstacles & 34.81 & \textbf{50.86} & \textbf{Reward-Based} \\
    Fast wind & 20.70 & 23.56 & Reward-Based \\
    Static environment & 51.01 & 56.47 & Reward-Based \\
    3 agents & 43.21 & 45.71 & Reward-Based \\
    10 agents & 17.15 & 23.13 & \textbf{Reward-Based} \\
    \bottomrule
  \end{tabular}
\end{table}

\subsubsection{Observations}

The experimental results reveal nuanced findings:

\begin{enumerate}
  \item \textbf{No universal winner}: First-Available outperforms Reward-Based in baseline, low-obstacle, fast-wind, and 3-agent scenarios.
  \item \textbf{Reward-Based excels in complex scenarios}: high obstacles (61\% faster completion, 46\% higher efficiency at 10 agents), static environment completes more holes before timeout.
  \item \textbf{Why First-Available sometimes wins}: lower overhead; in fast-changing environments Reward-Based may waste effort on relocated targets; with fewer agents the ``closest'' heuristic suffices.
  \item \textbf{Safety rate}: Both strategies show similar safety rates (0.64--0.78).
\end{enumerate}

\subsection{Parameter Effects Analysis}

\subsubsection{Effect of Obstacle Density}

\begin{table}[htbp]
  \centering
  \caption{Completion time vs obstacle count}
  \begin{tabular}{lcc}
    \toprule
    \sustechheaderrow
    \sustechtableheader{Obstacles} & \sustechtableheader{First-Avail. (ticks)} & \sustechtableheader{Reward-Based (ticks)} \\
    \midrule
    5 (low) & 844.5 & 538.3 \\
    15 (baseline) & 1309.1 & 1450.9 \\
    40 (high) & 974.1 & 376.0 \\
    \bottomrule
  \end{tabular}
\end{table}

\paragraph{Observations.}
\begin{itemize}
  \item \textbf{Low obstacles}: Reward-Based is 36\% faster (538 vs 844 ticks).
  \item \textbf{High obstacles}: Reward-Based is 61\% faster (376 vs 974 ticks).
  \item The gap widens as obstacle density increases---Reward-Based's path efficiency gains become more significant.
  \item First-Available is negatively affected by obstacles as it does not optimize path selection.
\end{itemize}

\subsubsection{Effect of Dynamic Update Interval}

\begin{table}[htbp]
  \centering
  \caption{Completion time vs wind interval}
  \begin{tabular}{lcc}
    \toprule
    \sustechheaderrow
    \sustechtableheader{Wind Interval} & \sustechtableheader{First-Avail.} & \sustechtableheader{Reward-Based} \\
    \midrule
    Static ($\infty$) & 5000 (4.0 holes left) & 5000 (2.4 holes left) \\
    Fast (10 ticks) & 1059.2 & 1424.3 \\
    \bottomrule
  \end{tabular}
\end{table}

\paragraph{Observations.}
\begin{itemize}
  \item In static environments, both strategies hit the 5000 tick limit, but Reward-Based completes more holes (2.4 remaining vs 4.0).
  \item In fast-changing environments, First-Available performs better (1059 vs 1424 ticks).
  \item This suggests Reward-Based may select targets that get relocated by wind, causing wasted effort.
  \item First-Available's ``grab first available'' approach is more robust to environmental changes.
\end{itemize}

\subsubsection{Effect of Agent Count (Non-Collaborative)}

\begin{table}[htbp]
  \centering
  \caption{Completion time vs agent count}
  \begin{tabular}{lcc}
    \toprule
    \sustechheaderrow
    \sustechtableheader{\# Agents} & \sustechtableheader{First-Avail.} & \sustechtableheader{Reward-Based} \\
    \midrule
    3 & 964.9 & 1364.9 \\
    5 (baseline) & 1309.1 & 1450.9 \\
    10 & 550.7 & 347.2 \\
    \bottomrule
  \end{tabular}
\end{table}

\paragraph{Observations.}
\begin{itemize}
  \item \textbf{With 3 agents}: First-Available is 41\% faster.
  \item \textbf{With 10 agents}: Reward-Based is 37\% faster.
  \item The crossover occurs because with more agents, competition for ``first available'' holes increases, while Reward-Based can assign diverse targets based on utility.
\end{itemize}

In non-collaborative mode, adding more agents does not always improve performance:
\begin{itemize}
  \item Initially improves as more holes can be worked on simultaneously.
  \item Eventually saturates due to limited number of holes/tiles.
  \item First-available shows more conflict (multiple agents targeting the same hole).
\end{itemize}

\subsection{Discussion}

\subsubsection{What Worked Well}

\begin{itemize}
  \item SEU-based selection balances reward, distance, and energy feasibility
  \item Proactive recharge prevents failures in challenging configurations
  \item Wind interval sensitivity is manageable with proper pathfinding
\end{itemize}

\subsubsection{Sources of Inefficiency}

\begin{itemize}
  \item No coordination: redundant or conflicting targets
  \item Static thresholds when far from stations
  \item Greedy pathfinding local minima under dense obstacles
\end{itemize}

\subsubsection{Potential Improvements}

\begin{itemize}
  \item Multi-agent coordination (R5)
  \item Dynamic threshold recharge strategy
  \item A* or D* Lite for complex obstacles
\end{itemize}

\subsection{Pathfinding Algorithm Comparison}
\label{sec:pathfinding-compare}

\subsubsection{Experimental Setup}

This section evaluates Greedy, A*, and D* Lite. Implementations come from \texttt{PathFinderAlgo-1.0-SNAPSHOT.jar} in \texttt{lib/}.

\begin{sustechnote}[title={PathFinderAlgo}]
Java pathfinding library (A*, D* Lite, variants). Repository: \url{https://github.com/Wellshh/PathFinderAlgo}. This project adds grid adapters, sensing, and toroidal borders.
\end{sustechnote}

\begin{table}[htbp]
  \centering
  \caption{Pathfinding test configuration}
  \begin{tabular}{lc}
    \toprule
    \sustechheaderrow
    \sustechtableheader{Parameter} & \sustechtableheader{Value} \\
    \midrule
    Grid Size & $50{\times}50$ \\
    Agents & 5 \\
    Holes & 10 \\
    Obstacles & 15 \\
    Wind Interval & 20 ticks (dynamic) \\
    Pathfinding & Greedy / A* / D* Lite \\
    \bottomrule
  \end{tabular}
\end{table}

\subsubsection{Results}

\begin{table}[htbp]
  \centering
  \caption{Pathfinding algorithm performance}
  \begin{tabular}{lcccc}
    \toprule
    \sustechheaderrow
    \sustechtableheader{Algorithm} & \sustechtableheader{Time (ticks)} & \sustechtableheader{Reward Eff.} & \sustechtableheader{Path Cost} & \sustechtableheader{Safety} \\
    \midrule
    \textbf{Greedy} & 2084.8 & 8.97 & 2.55 & 100\% \\
    \textbf{A$^\ast$} & 2998.0 & 7.35 & 2.16 & 58\% \\
    \textbf{D$^\ast$\,Lite} & 4298.2 & 10.38 & 1.78 & 68\% \\
    \bottomrule
  \end{tabular}
\end{table}

\textit{Lower completion time and path cost are better. Higher reward efficiency and safety rate are better.}

\tileworldfigure{Multi-metric comparison of pathfinding algorithms (Greedy, A*, D* Lite).}{assets/pathfinding_comparison.png}{fig:pathfinding}

\textbf{Conclusion}: No universal winner---use Greedy for speed, D* Lite for efficiency; A* is risky under frequent full replans in dynamic wind.

\section{R5: Multi-Agent Collaboration}

\subsection{Collaboration Mechanism: Contract Net Protocol}

The multi-agent collaboration uses the \textbf{Contract Net Protocol (CNP)} (Smith, 1980), a market-based mechanism without centralized control.

\begin{sustechnote}[title={References}]
Smith, R.G. (1980). \emph{The Contract Net Protocol}. \url{https://www.reidgsmith.com/The_Contract_Net_Protocol_Dec-1980.pdf}. Wikipedia: \url{https://en.wikipedia.org/wiki/Contract_Net_Protocol}
\end{sustechnote}

\subsubsection{CNP Protocol Overview}

The Contract Net Protocol operates through a four-phase handshake between a \emph{manager} (the agent that announces a task) and \emph{contractors} (peers that may satisfy it). \Cref{fig:cnp-handshake} summarizes the message flow; the following steps map this abstraction onto the simulation's \texttt{ContractNetBoard}.

\begin{figure}[htbp]
  \centering
  \begin{sustechsystemdiagram}
    % Tall lifeline-style entities (wider + taller than default sustech module)
    \sustechmodulebox[minimum width=36mm, minimum height=56mm, inner sep=7pt]{mgr}{at={(0,0)}}{%
      \textbf{Manager}\\[3pt]\scriptsize (Agent)}
    \sustechmodulebox[minimum width=36mm, minimum height=56mm, inner sep=7pt]{ctr}{right=78mm of mgr}{%
      \textbf{Contractors}\\[3pt]\scriptsize (Agents)}
    % Message attachment points along full vertical extent of each column
    \path (mgr.north east) -- (mgr.south east)
      coordinate[pos=0.10] (lane-m-1)
      coordinate[pos=0.36] (lane-m-2)
      coordinate[pos=0.64] (lane-m-3)
      coordinate[pos=0.90] (lane-m-4);
    \path (ctr.north west) -- (ctr.south west)
      coordinate[pos=0.10] (lane-c-1)
      coordinate[pos=0.36] (lane-c-2)
      coordinate[pos=0.64] (lane-c-3)
      coordinate[pos=0.90] (lane-c-4);
    \sustechconnecttowith{sustech highlight arrow}[][1.\ Call for proposals (CFP)]{lane-m-1}{lane-c-1}
    \sustechconnecttowith{sustech flow arrow}[][2.\ Submit bids]{lane-c-2}{lane-m-2}
    \sustechconnecttowith{sustech flow arrow}[][3.\ Award contract]{lane-m-3}{lane-c-3}
    \draw[sustech flow arrow, {Latex[length=2.5mm,width=1.8mm]}-{Latex[length=2.5mm,width=1.8mm]}]
      (lane-m-4) -- node[sustech edge label]{4.\ Execute / report} (lane-c-4);
  \end{sustechsystemdiagram}
  \caption{Contract Net handshake (conceptual)}
  \label{fig:cnp-handshake}
\end{figure}

\begin{enumerate}
  \item \textbf{Call For Proposals (CFP):} When an agent discovers an unclaimed task (a hole or an energy station), it announces the task to the \texttt{ContractNetBoard}, which broadcasts a call for proposals to eligible peers.
  \item \textbf{Bidding:} Interested agents respond with bids that include utility scores reflecting expected reward, travel cost, energy risk, and (where applicable) switching penalties---so the board can compare commitments on a common scale.
  \item \textbf{Award:} The board evaluates feasible bids and awards the contract to the highest-utility bidder; ties or infeasible bids are resolved according to the configured evaluation strategy.
  \item \textbf{Execution:} The awarded agent executes the task (acquire/move/fill or recharge). If execution fails or the task becomes invalid, the contract may be released and the task re-auctioned in a later tick.
\end{enumerate}

\subsubsection{Task Types}

The implementation classifies work into three CNP task types so the board can announce the minimal commitment required---from a full hole-servicing chain down to a direct fill or a recharge slot.

\begin{table}[htbp]
  \centering
  \caption{CNP task types}
  \begin{tabular}{lp{0.40\linewidth}p{0.34\linewidth}}
    \toprule
    \sustechheaderrow
    \sustechtableheader{Task Type} & \sustechtableheader{Description} & \sustechtableheader{When Announced} \\
    \midrule
    \texttt{SERVICE\_HOLE} &
      Complete task chain: acquire a tile, reach the hole, and fill it. &
      Agent still needs both a suitable tile and a target hole (full service contract). \\
    \texttt{DIRECT\_FILL\_HOLE} &
      Fill the hole only; the agent already carries a tile. &
      Agent holds a tile and only needs routing and execution for the fill action. \\
    \texttt{RECHARGE\_SLOT} &
      Claim a slot at an energy station for replenishment. &
      Agent must recharge (proactive or critical energy state) and competes for station access. \\
    \bottomrule
  \end{tabular}
\end{table}

Each row corresponds to a distinct announcement path on the board: \texttt{SERVICE\_HOLE} coordinates tile acquisition and hole completion, \texttt{DIRECT\_FILL\_HOLE} avoids redundant tile claims when cargo is already held, and \texttt{RECHARGE\_SLOT} serializes contention at scarce stations alongside hole tasks.

\subsection{Contract Evaluation Strategy}

The core bid evaluation is implemented in \texttt{ContractEvaluationStrategy.java}:

\begin{sustechcode}{Java}
public interface ContractEvaluationStrategy {
    ContractBid buildHoleBid(Agent bidder, Hole hole, List<Tile> candidateTiles,
            Grid<Object> grid, List<EnergyStation> allStations,
            int safetyMargin, int maxGridDistance);

    ContractBid buildRechargeBid(Agent bidder, EnergyStation station,
            Grid<Object> grid, int safetyMargin, int maxGridDistance);

    int compareBids(ContractBid a, ContractBid b);
}
\end{sustechcode}

Utility formulas for hole service, direct fill, and recharge follow the same weighted structure as in R4 (score, hold bonus, distance, risk, switch penalty; recharge adds urgency and queue terms). Feasibility gates require available targets, sufficient energy, and agent not powered down.

\subsection{Integration with BDI}

In Contract Net mode, coordination goes through the board each tick:

\begin{sustechcode}{Java}
if (this.contractNetBoard != null) {
    this.contractNetBoard.tick(this, beliefState, currentPosition);
} else {
    List<Agent> partners = this.policyBundle.getCoordinationStrategy()
            .selectPartners(this, beliefState, this.currentIntention);
}
\end{sustechcode}

\subsection{Evaluation: Collaboration vs Non-Collaboration}

\begin{table}[htbp]
  \centering
  \caption{Coordination mode comparison (seeds 1--10, A* pathfinding)}
  \begin{tabular}{lcccc}
    \toprule
    \sustechheaderrow
    \sustechtableheader{Metric} & \sustechtableheader{DISABLED} & \sustechtableheader{PASSIVE} & \sustechtableheader{ACTIVE} & \sustechtableheader{CONTRACT\_NET} \\
    \midrule
    Completion Time & 2998.0 & 3063.1 & 3278.0 & \textbf{699.8} * \\
    Reward Efficiency & 7.71 & 7.54 & 7.52 & \textbf{19.48} * \\
    Path Cost & 2.14 & 2.14 & 2.17 & \textbf{1.51} * \\
    Safety Rate & 58\% & 58\% & 56\% & \textbf{100\%} * \\
    Remaining Holes & 0.6 & 0.6 & 1.1 & \textbf{0.0} * \\
    Cooperation Gain & 0.0 & -2.17 & -9.34 & \textbf{+76.66} * \\
    \bottomrule
  \end{tabular}
\end{table}

\tileworldfigure{Radar chart comparing coordination modes across metrics. CONTRACT\_NET dominates.}{assets/coordination_radar.png}{fig:coord-radar}

\subsubsection{Key Observations}

Contract Net yields about $4.3\times$ faster completion, $\sim 3\times$ reward efficiency, 100\% safety, and zero remaining holes versus other modes; ACTIVE adds overhead without effective allocation.

\subsection{Discussion}

CNP improves efficiency by assigning high-value holes to capable nearby agents, prioritizing critical recharge, and reducing redundant work. Limitations: single-agent task chains, CNP latency under very fast wind, static bid weights.

\begin{thebibliography}{9}
\bibitem{pollack1990}
M.~E. Pollack and M.~Ringuette.
\newblock Introducing the Tileworld: Experimentally Evaluating Agent Architectures.
\newblock In \emph{AAAI-90}, pages 183--189. AAAI Press, 1990.

\bibitem{smith1980}
R.~G. Smith.
\newblock The Contract Net Protocol: High-Level Communication and Control in a Distributed Problem Solver.
\newblock \emph{IEEE Trans.\ Comput.}, C-29(12):1104--1113, 1980.
\newblock \url{https://www.reidgsmith.com/The_Contract_Net_Protocol_Dec-1980.pdf}

\bibitem{wellshh2024}
Wellshh.
\newblock \emph{PathFinderAlgo --- Path Planning Algorithms Library}, 2024.
\newblock \url{https://github.com/Wellshh/PathFinderAlgo}

\bibitem{wooldridge2003}
M.~Wooldridge.
\newblock \emph{Reasoning about Rational Agents}.
\newblock MIT Press, 2003.
\end{thebibliography}

\appendix
\section{Key Code Snippets}

\subsection{BDI Execution Cycle (\texttt{Agent.java\#run()})}

\begin{sustechcode}{Java}
@ScheduledMethod(start = 1, interval = 1)
public void run() {
    if (this.poweredDown) { return; }

    AgentBeliefState preBelief = this.buildBeliefState();
    if (this.contractNetBoard == null) {
        this.reviseBeliefs(preBelief);
    }
    AgentBeliefState beliefState = this.buildBeliefState();

    List<AgentDesire> desires = this.policyBundle.getDesireEvaluationStrategy()
            .evaluate(beliefState, this);

    AgentIntention nextIntention = this.policyBundle.getTaskSelectionStrategy()
            .selectIntention(beliefState, desires, this.currentIntention);
    this.commitIntention(nextIntention, currentPosition);

    List<Agent> partners = this.policyBundle.getCoordinationStrategy()
            .selectPartners(this, beliefState, this.currentIntention);
    this.syncCoordinationPartners(partners);

    this.act(context);
}
\end{sustechcode}

\subsection{Desire Evaluation (\texttt{DefaultDesireEvaluationStrategy.java})}

\begin{sustechcode}{Java}
public List<AgentDesire> evaluate(AgentBeliefState beliefState, Agent agent) {
    List<AgentDesire> desires = new ArrayList<>();

    if (this.rechargeStrategy.shouldRecharge(beliefState)
            && beliefState.getClaimedEnergyStation() != null) {
        desires.add(new AgentDesire(AgentDesireType.RECHARGE, 100,
                beliefState.getClaimedEnergyStation()));
    }

    if (beliefState.isHoldingTile() && beliefState.getClaimedHole() != null) {
        desires.add(new AgentDesire(AgentDesireType.FILL_HOLE, 90,
                beliefState.getClaimedHole()));
    }

    if (!beliefState.isHoldingTile() && beliefState.getClaimedTile() != null) {
        desires.add(new AgentDesire(AgentDesireType.ACQUIRE_TILE, 80,
                beliefState.getClaimedTile()));
    }

    desires.add(new AgentDesire(AgentDesireType.EXPLORE, 10, null));
    return desires;
}
\end{sustechcode}

\subsection{Greedy Path Planning (\texttt{GreedyPathPlanningStrategy.java})}

\begin{sustechcode}{Java}
public GridPoint nextStep(Agent agent, GridPoint currentPosition,
        GridPoint targetPosition) {
    GridDimensions dimensions = agent.getGrid().getDimensions();
    int curX = currentPosition.getX(), curY = currentPosition.getY();

    if (curX != targetPosition.getX()) {
        int[] candidateX = this.computeNextCoordinate(
                curX, targetPosition.getX(), dimensions.getWidth());
        if (agent.canMoveTo(candidateX[0], curY))
            return new GridPoint(candidateX[0], curY);
        if (agent.canMoveTo(candidateX[1], curY))
            return new GridPoint(candidateX[1], curY);
    }

    if (curY != targetPosition.getY()) {
        int[] candidateY = this.computeNextCoordinate(
                curY, targetPosition.getY(), dimensions.getHeight());
        if (agent.canMoveTo(curX, candidateY[0]))
            return new GridPoint(curX, candidateY[0]);
        if (agent.canMoveTo(curX, candidateY[1]))
            return new GridPoint(curX, candidateY[1]);
    }

    return currentPosition;
}
\end{sustechcode}

\subsection{Wind Dynamic Update (\texttt{Wind.java\#blow()})}

\begin{sustechcode}{Java}
public void blow() {
    List<Obstacle> obstacles = new ArrayList<>();
    List<Tile> tiles = new ArrayList<>();
    for (Object obj : this.context) {
        if (obj instanceof Obstacle) obstacles.add((Obstacle) obj);
        else if (obj instanceof Tile) tiles.add((Tile) obj);
    }

    List<int[]> candidateCells = this.collectCandidateCells();
    this.pickWithoutReplacement(candidateCells, obstacles.size() + tiles.size());

    int cursor = 0;
    for (Obstacle o : obstacles) {
        int[] pos = candidateCells.get(cursor++);
        this.grid.moveTo(o, pos[0], pos[1]);
    }
    for (Tile t : tiles) {
        int[] pos = candidateCells.get(cursor++);
        this.grid.moveTo(t, pos[0], pos[1]);
    }
}
\end{sustechcode}
